% =====================================================================
%  Chapter 19 — The Proof in Natural Language, 3
% =====================================================================

\chapter{The Arithmetic Factor}
\label{ch:count}

This chapter is the heart of the method, and the place where it is most
unlike the function-field proof it descends from.

\section{The test-function sum}Fix a smooth compactly supported even function $t$ on $\mathbb{R}$, and fix
the integer $k$. Consider
\begin{equation}
S(t,k) \;=\; \sum_{\substack{h \bmod q \\ x_1,\dots,x_k \bmod q}}
t(x_1)\cdots t(x_k)\;e_q\!\left(h\bigl(x_1x_2 + x_2x_3 + \cdots + x_{k-1}x_k\bigr)\right),
\label{eq:tricksum}
\end{equation}
where $e_q$ is the standard additive character of $\mathbb{F}_q$ and the sum
runs over the coefficients $h$ of a quadratic form in $k$ variables together
with the $k$ coordinates.

The remark's item 5 states this sum, and the instruction that makes it work:
\emph{sum over $h$ first}.

\section{Why the order of summation matters}For fixed $x_1,\dots,x_k$, the inner sum over $h$ is a character sum over the
linear functional $x_1x_2 + \cdots + x_{k-1}x_k$. It is either zero --- when
the functional is zero --- or has absolute value $q^{1/2}$. So the inner sum
does not vanish: it is a Gauss sum.

Doing the $h$-sum first replaces a sum over a group of characters by a sum
over the solutions of a quadratic equation. That is a \emph{point count}, and
a point count is arithmetic: it is bounded, in the function-field setting, by
the Weil bound, and its excess over the main term is the square-root-sized
quantity the argument needs.

\begin{quote}
\textbf{Proposition E.} The inner sum over $h$ in \eqref{eq:tricksum} is a
Gauss sum of absolute value $q^{1/2}$, except where the quadratic form is
degenerate on the tuple. Consequently the total is a sum of Gauss sums
weighted by $t(x_1)\cdots t(x_k)$, and the arithmetic of the problem enters
through the count of the tuples for which the form is non-degenerate.
\end{quote}

The formal counterpart is the local additive character expansion in
\code{Formalism/SharpRestriction/Divisibility/RankTruncation/LocalAdditiveCharacterExpansion.lean},
with the character sum itself in
\code{.../Divisibility/CutDivD.lean} as \code{sum\_char}.

\section{The sign from the middle cohomology}The next move is the one that makes the argument work over the integers.

Consider a quadric $Q$ in $k$ variables: the zero locus of
$x_1x_2 + \cdots + x_{k-1}x_k$. The count of its $\mathbb{F}_q$-points, and
more particularly the character sum over it, is determined --- up to a sign
and a known factor --- by the middle cohomology of $Q$.

\begin{quote}
\textbf{Proposition F.} For a non-degenerate quadric in $k$ variables over
$\mathbb{F}_q$, the Gauss sum attached to its defining form is, up to a sign
that is fixed by the parity of $k$ and by the rank of the form, equal to a
power of $q^{1/2}$ times a product of local Gauss sums.
\end{quote}

The remark's phrase is that ``Bertini gives a definite sign from middle
cohomology''. Bertini's theorem is what identifies the middle cohomology of a
smooth quadric with the exterior algebra on the Weil representation, and it
is why the sign of the sum is computable rather than conjectural.

Here is the crucial point for the method. Over $\mathbb{F}_q$ this argument
supplies a sign because the quadric is a variety over a finite field and its
cohomology is that of a variety. Over $\mathbb{Z}$ there is no such thing ---
but the formalisation asserts that the same construction, carried out in the
automorphic setting, produces the same sign. That is the content of the
``natively-automorphic'' claim, and it is the claim with the fewest classical
precedents.

The formal counterpart is the regularised Siegel--Weil development,
\code{Development/WeilBound/RepresentationTheoryGL/RegularisedSiegelWeil/},
153 files, and the quadric-related statements in
\code{Development/WeilBound/ThetaLifts/LeviDeterminants/}.

\section{The decomposition \texorpdfstring{$qN = Mk + R$}{qN = Mk + R}}Let $N$ be the number in question, and let the sum be split into its trivial
part and its excess. The remark's item 3 gives the shape:

\begin{quote}
\small
``$qN = Mk + R$, using $x_1\cdot x_2 + \cdots + x_{k-1}\cdot x_k$, wanting
`automorphy' of $N$ since $\mu$ is automorphic from $\theta$ on sixfold cover
of $\mathrm{GL}_5$ as a period integral of thetas.''
\end{quote}

\begin{quote}
\textbf{Proposition G.} The count splits as
\[
qN = Mk + R,
\]
where $M = M_k$ is the main term --- the number of trivial solutions of the
quadratic equation, a quantity depending only on $k$ --- and $R$ is the
remainder.
\end{quote}

The main term is completely understood: for a quadratic form in $k$ variables
the number of trivial solutions is a fixed number depending only on $k$ (for
the hyperbolic form, $q^{k-1} + O(q^{k/2})$). The information is entirely in
$R$.

\section{The square-root requirement}\begin{quote}
\textbf{Proposition H.} For the argument to yield the Weil bound, $R$ must
satisfy
\[
R = O\!\left(q^{k/2+\varepsilon}\right),
\]
for every $\varepsilon > 0$, with the constant controlled.
\end{quote}

Why $k/2$? Because the Weil bound on a curve of genus $g$ over $\mathbb{F}_q$ is
$\#C(\mathbb{F}_q) \le q + 1 + O(g\sqrt q)$, and the genus is of order
$q^{k/2}$ when the geometry is that of a quadric in $k$ variables. So the
requirement on $R$ is exactly the requirement that the excess over the main
term be of the size of the Weil bound for the geometry in question.

This is where the remark's phrase ``$R$ square-root-quality up to an inch or
so as Katz would say'' sits. The phrase is a hedge: the bound is
$q^{k/2}$ up to a bounded loss, not exactly $q^{k/2}$, and the author is
careful to say so.

\section{Why the square root is hard, and how it is obtained}The Weil bound is not free: it is a theorem about the eigenvalues of Frobenius
acting on the cohomology of the variety. If the variety is a curve, its
cohomology is small and the eigenvalues are controlled. In the present
construction the geometry is that of a quadric in $k$ variables, whose
middle cohomology is large, and the eigenvalues that must be controlled are
those of the \emph{automorphic} correspondence --- which is a much harder
object to control than a finite-field Frobenius.

This is where the two ingredients the remark lists come in.

\begin{description}[style=nextline,leftmargin=0pt,labelwidth=0pt,labelsep=0pt,font=\sffamily\bfseries\small,
  itemsep=6pt]
\item[A main term.] To get the square-root bound one needs the leading term
of the relevant $L$-function to be known. The remark's item 14 lists this as
the first of three attempts: ``attempt to get a main term''.
\item[Parseval, iterated $k$ times.] The bound on the excess requires the
contribution of each sign class to be separated by root. The remark obtains
this by iterating the Parseval identity $k$ times.
\end{description}

The formalisation of the Parseval step is
\code{ConditionalKappaBounds.L2.LocalAdditiveCharacterExpansion} in
\code{.../Divisibility/RankTruncation/LocalAdditiveCharacterExpansion.lean},
and the separation condition it produces is the cubic gap of
Chapter~\ref{ch:step3} --- the requirement that distinct roots be separated
by the cube of the weight-and-level combination.

\section{Where the square root is still open}It should be said plainly: the proposition that the excess $R$ has
square-root quality, for the automorphic construction over the integers, is
not a theorem anyone has published. The Weil bound for curves is a theorem; its
automorphic counterpart in this generality is the thing the method is
proposing, and it is exactly the thing the remark's three attempts are aimed
at.

What the formalisation does supply is the machinery: the character expansion,
the local bounds, the separation condition, and the budget arithmetic of
Chapter~\ref{ch:budget} that decides whether the conditions can be met
simultaneously.

\begin{keytakeaways}
\item The test-function sum \eqref{eq:tricksum}, summed over $h$ first,
turns the problem into a point count, which is arithmetic.
\item The sign comes from the middle cohomology of the quadric, by Bertini, and
is the point at which the argument departs from the finite-field setting.
\item The count splits as $qN = Mk + R$, with the main term understood and
all the information in the remainder.
\item The requirement is that $R$ have square-root quality, and obtaining it
automorphically is the new content of the method --- and it is not published.
\end{keytakeaways}
